% topology: geometry — 物质坐标映射 [0,R(t)] -> [0,1]
% MH-TIKZ-ABSOLUTE-OK: 坐标对应物理域与计算域的位置
\documentclass[border=4pt]{standalone}
\usepackage{ctex}
\usepackage{tikz}
\usepackage{amsmath}
\usetikzlibrary{arrows.meta,calc,decorations.pathreplacing}
\definecolor{paperInk}{RGB}{0,0,0}
\definecolor{paperAccent}{RGB}{58,84,156}
\begin{document}
\begin{tikzpicture}[
  >={Stealth[length=5pt]},
  axisline/.style={paperInk,line width=.9pt,-{Stealth[length=5pt]}},
  move/.style={paperAccent,line width=1pt,-{Stealth[length=4pt]}},
  tick/.style={paperInk,line width=.8pt}]

  % ===== 左：物理域 [0, R(t)]，边界随时间内移 =====
  \def\Lx{0}
  \draw[axisline] (\Lx-0.3,0) -- (\Lx+4.6,0) node[right]{$r$};
  \node[paperInk,tick] at (\Lx,0) [below=3pt]{$0$};
  \draw[tick] (\Lx,-0.12)--(\Lx,0.12);
  % 初始边界 R0
  \draw[tick] (\Lx+4.0,-0.12)--(\Lx+4.0,0.12);
  \node[paperInk] at (\Lx+4.0,0)[below=3pt]{$R_0$};
  % 当前收缩边界 R(t)
  \draw[paperAccent,line width=1.1pt] (\Lx+2.6,-0.35)--(\Lx+2.6,0.35);
  \node[paperAccent] at (\Lx+2.6,0)[above=4pt]{$R(t)$};
  % 边界内移箭头 + 固相迁移速度
  \draw[move] (\Lx+3.9,0.7) -- (\Lx+2.7,0.7)
      node[midway,above,font=\small]{边界内移};
  \foreach \x in {1.0,1.8,2.3}{
    \draw[move] (\Lx+\x+0.35,-0.62)--(\Lx+\x-0.05,-0.62);
  }
  \node[paperAccent,font=\small,align=center] at (\Lx+2.0,-1.05)
      {固相迁移 $u_s=\dfrac{r}{R}\dot R$};
  \node[paperInk] at (\Lx+2.0,1.5){物理域（动边界）};

  % ===== 映射箭头 =====
  \draw[paperAccent,line width=1.3pt,-{Stealth[length=6pt]}]
      (5.0,0) .. controls (5.8,0.5) and (6.2,0.5) .. (7.0,0);
  \node[paperAccent,align=center,font=\small] at (6.0,0.95)
      {$\eta=\dfrac{r}{R(t)}$};

  % ===== 右：计算域 [0,1] 固定 =====
  \def\Rx{7.4}
  \draw[axisline] (\Rx-0.3,0) -- (\Rx+4.2,0) node[right]{$\eta$};
  \draw[tick] (\Rx,-0.12)--(\Rx,0.12);
  \node[paperInk] at (\Rx,0)[below=3pt]{$0$};
  \draw[tick] (\Rx+3.6,-0.12)--(\Rx+3.6,0.12);
  \node[paperInk] at (\Rx+3.6,0)[below=3pt]{$1$};
  % 固定网格（等分刻度）
  \foreach \k in {1,...,5}{
    \draw[paperInk,line width=.5pt] (\Rx+3.6*\k/6,-0.09)--(\Rx+3.6*\k/6,0.09);
  }
  \node[paperInk] at (\Rx+1.8,1.5){计算域（固定网格）};
  \node[paperInk,font=\small] at (\Rx+1.8,-1.05){域固定，收缩改由系数承担};

  % ===== 底部：相消结论 =====
  \node[paperAccent,align=center,font=\small] at (5.9,-2.2)
    {坐标运动项 $-\dfrac{\eta\dot R}{R}\partial_\eta C$ 与固相对流项 $+\dfrac{\eta\dot R}{R}\partial_\eta C$ 精确相消};
  \node[paperInk,align=center,font=\small] at (5.9,-3.0)
    {收缩效应仅以因子 $\dfrac{1}{R^2(t)}$ 进入扩散项，方程不含伪对流项};
\end{tikzpicture}
\end{document}
